% ============================================================================
% Section 8.4 — Comparing Populations with Measures (AK/OK/TX: extended measures emphasis)
% CCSS 7.SP.B.4
% ============================================================================
\section{Comparing Populations with Measures of Center and Variability}

\begin{stepGoal}
\begin{itemize}
    \item Use mean, median, range, and IQR to compare two data sets
    \item Explain why IQR is often better than range for comparison
\end{itemize}
\end{stepGoal}

\begin{stepVocabBox}
    \vocabItem{Mean}{Sum of all values divided by the count.}
    \vocabItem{Median}{The middle value when data is in order.}
    \vocabItem{Range}{Highest $-$ lowest — shows the total spread.}
    \vocabItem{Interquartile Range (IQR)}{$Q_3 - Q_1$ — the spread of the \textbf{middle $50\%$} of the data.}
\end{stepVocabBox}

\begin{stepsCard}[How to Compare Two Populations with Extended Measures]
    \stepItem{Order each data set and find the \textbf{mean} and \textbf{median}.}
    \stepItem{Find the \textbf{range} (highest $-$ lowest).}
    \stepItem{Find $Q_1$, $Q_3$, and the \textbf{IQR} ($Q_3 - Q_1$).}
    \stepItem{\textbf{Compare}: Which group has a higher center? Which has less spread?}
    \stepItem{Write a \textbf{conclusion} supported by at least two measures.}
\end{stepsCard}

% --- Worked Example ---
\begin{stepExample}{Crew A daily catch (lb): $120, 135, 140, 150, 155$. Crew B: $90, 130, 145, 160, 200$. Compare.}
    \stepShow{1} Crew A mean $= \frac{700}{5} = 140$, median $= 140$.\\
    Crew B mean $= \frac{725}{5} = 145$, median $= 145$.

    \stepShow{2} Crew A range $= 155 - 120 = 35$. Crew B range $= 200 - 90 = 110$.

    \stepShow{3} Crew A: $Q_1 = 135$, $Q_3 = 150$, IQR $= 15$.\\
    Crew B: $Q_1 = 130$, $Q_3 = 160$, IQR $= 30$.

    \stepShow{4} Crew B averages slightly more ($145$ vs.\ $140$), but Crew A is far more consistent (IQR $15$ vs.\ $30$).

    \stepShow{5} Crew B's range ($110$) is inflated by one large catch. IQR gives a more reliable picture.
    \stepResult{Crew B averages more, but Crew A is much more consistent (IQR $15$ vs.\ $30$).}
\end{stepExample}

\watchOut{Range can be thrown off by a single extreme value. Always check IQR for a more stable comparison!}

% ============================================================================
% PRACTICE
% ============================================================================
\newpage
\begin{stepPractice}{Comparing Populations with Measures Practice}
    \resetProblems

    \stepPracticeHeader[step1Color]{Find Center}
    \prob Find the mean and median: $5, 8, 10, 12, 15$. \answerBlank[3cm]
    \answer{Mean $= 10$, median $= 10$}

    \stepPracticeHeader[step2Color]{Find Spread}
    \prob Find the range and IQR of $5, 8, 10, 12, 15$. \answerBlank[3cm]
    \answer{Range $= 10$; $Q_1 = 8$, $Q_3 = 12$, IQR $= 4$}

    \stepPracticeHeader[step3Color]{Why IQR?}
    \prob Data: $5, 8, 10, 12, 15, 50$. Range $= 45$ and IQR $= 7$. Which better shows typical spread? \answerBlank[3cm]
    \answer{IQR — the outlier $50$ inflates the range; IQR ignores it}

    \stepPracticeHeader[step4Color]{Compare Two Groups}
    \prob Group A: mean $72$, IQR $= 6$. Group B: mean $78$, IQR $= 20$. Which is more consistent? \answerBlank[3cm]
    \answer{Group A — smaller IQR means less variability}

    \wordProblem{Student X daily steps: $7{,}000$; $8{,}000$; $8{,}500$; $9{,}000$; $9{,}500$. Student Y: $3{,}000$; $7{,}500$; $8{,}000$; $10{,}000$; $13{,}500$. Compare using mean, range, and IQR.}{~}
    \answer{X: mean $= 8{,}400$, range $= 2{,}500$, IQR $= 1{,}000$. Y: mean $= 8{,}400$, range $= 10{,}500$, IQR $= 2{,}500$. Same average, but X is far more consistent.}
\end{stepPractice}
